\begin{frame}{What students must learn to handle}

\begin{alertblock}{\faBolt~What becomes visible}
The design makes certain challenges explicit
\end{alertblock}

\begin{itemize}
    \item giving and receiving critical feedback
    \item making work visible and reviewable
    \item managing workload in a structured workflow
\end{itemize}
\end{frame}
\note{
What becomes clear with this design  
is not just improved outcomes --  
but also the challenges students need to learn to handle.  

\begin{enumerate}
    \item \emph{Giving and receiving critical feedback}:  
    Students are often hesitant to critique peers,  
    but this becomes a central part of the workflow.  

    \item \emph{Making work visible}:  
    Everything is open to review --  
    which creates accountability,  
    but also requires students to articulate and justify their decisions.  
    
    \item \emph{Managing workload}:  
    The structure introduces overhead,  
    and students feel that early on.  
    
    But that structure is also what enables  
    iteration, feedback, and improvement.  
    
\end{enumerate}


So the point is not to remove these challenges.  

They are a direct consequence of making the work visible --  
and part of what students need to learn to navigate.
}